\documentclass[10pt]{article}
\usepackage[T1]{fontenc}
\usepackage[utf8]{inputenc}
\usepackage[french]{babel}
\usepackage{amsmath,amssymb,mathrsfs,mathpazo}
\usepackage{geometry}
\geometry{a4paper,margin=2cm}
\usepackage{graphicx,array,tabularx,tabularray,colortbl,multirow,multicol,booktabs,enumitem,pifont,ulem,fancybox,diagbox,cancel}
\usepackage{pstricks,pstricks-add,pst-plot,pst-tree,pst-node,pst-text,pst-func,pst-math,pst-bspline,pst-3dplot,pst-eucl,pst-solides3d}
\let\pscalculate\undefined
\usepackage{pst-bezier}
\usepackage{tikz,pgfplots,tkz-tab}
\usetikzlibrary{arrows,arrows.meta,patterns,calc,positioning,trees,decorations.markings,shapes.geometric,shapes.misc,shapes.arrows,intersections,decorations.pathmorphing,decorations.pathreplacing,backgrounds,shadows,babel}
\pgfplotsset{compat=1.18}
\usepgfplotslibrary{fillbetween}
\usepackage[np]{numprint}
\usepackage{listings,makecell,pas-tableur,twoopt}
\let\barre\undefined
\usepackage[pstricks]{tablvar}
\definecolor{darkgreen}{RGB}{0,100,0}
\providecommand{\red}{\color{red}}
\providecommand{\blue}{\color{blue}}
\providecommand{\green}{\color{green}}
\providecommand{\black}{\color{black}}
\providecommand{\white}{\color{white}}
\providecommand{\gray}{\color{gray}}
\providecommand{\cyan}{\color{cyan}}
\providecommand{\magenta}{\color{magenta}}
\providecommand{\yellow}{\color{yellow}}
\providecommand{\violet}{\color{violet}}
\providecommand{\orange}{\color{orange}}
\providecommand{\euro}{\texteuro}
\providecommand{\EUR}{\texteuro}
\providecommand{\ieme}{\textsuperscript{e}}
\providecommand{\degres}{\ensuremath{^\circ}}
\providecommand{\degree}{\ensuremath{^\circ}}
\providecommand{\vv}[1]{\overrightarrow{#1}}
\providecommand{\Coord}[1]{#1}
\providecommand{\decosix}{}
\providecommand{\starredbullet}{\ensuremath{\bullet}}
\pagestyle{empty}
\setlength{\parindent}{0pt}
\usepackage{adjustbox}\DeclareUnicodeCharacter{2236}{:}\DeclareUnicodeCharacter{2190}{\ensuremath{\leftarrow}}
\providecommand{\R}{\mathbb{R}}
\providecommand{\N}{\mathbb{N}}
\providecommand{\Z}{\mathbb{Z}}
\providecommand{\Q}{\mathbb{Q}}
\providecommand{\D}{\mathbb{D}}
\providecommand{\C}{\mathbb{C}}
\providecommand{\e}{{\mathrm{e}}}
\providecommand{\vect}[1]{\overrightarrow{#1}}
\providecommand{\vectt}[1]{\overrightarrow{\mathrm{#1}}}
\providecommand{\vvt}[1]{\overrightarrow{\mathrm{#1}}}
\providecommand{\np}[2][]{{#2~#1}}
\providecommand{\vv}[1]{{\overrightarrow{#1}}}
\renewcommand{\d}{{\mathrm{d}}}
\providecommand{\degres}{{^\circ}}
\providecommand{\degre}{{^\circ}}
\providecommand{\blue}{\color{blue}}
\providecommand{\red}{\color{red}}
\providecommand{\enskip}{\quad}
\providecommand{\llbracket}{\lbrack\!\lbrack}
\providecommand{\rrbracket}{\rbrack\!\rbrack}
\providecommand{\euro}{{\text{€}}}
\providecommand{\no}{{\text{n°}}}
\providecommand{\interval}[2]{{[#1;#2]}}
\providecommand{\intervFO}[2]{{[#1;#2[}}
\providecommand{\dotfill}{\ldots}

\providecommand{\ds}{\displaystyle}
\providecommand{\textding}[1]{#1}
\providecommand{\Oij}{\(\mathrm{O};\vec{\imath},\vec{\jmath}\)}
\providecommand{\Oijk}{\(\mathrm{O};\vec{\imath},\vec{\jmath},\vec{k}\)}
\providecommand{\Ouv}{\(\mathrm{O};\vec{u},\vec{v}\)}
\providecommand{\intervFF}[2]{\ensuremath{[#1;#2]}}
\renewcommand{\euro}{\ifmmode\text{\texteuro}\else\texteuro\fi}
\DeclareUnicodeCharacter{00B0}{\ifmmode^{\circ}\else\textdegree\fi}
\DeclareUnicodeCharacter{20AC}{\ifmmode\text{\texteuro}\else\texteuro\fi}
% Énoncés : Éducation nationale. Transcription et corrigés : APMEP, Denis Vergès et contributeurs. Découpage et classement : CoopMaths, Sébastien Lozano, Éric Elter et contributeurs. Adaptation de la présentation : Maths Explorer.
% CC BY-SA : https://coopmaths.fr/www/about/
% https://coopmaths.fr/alea/static/eam/2026/tex/eam_specifique_2026_09_metropole_3.tex
\begin{document}
\noindent\textit{APMEP / CoopMaths. CC BY-SA. Présentation adaptée par Maths Explorer.}\par\medskip
Une étude a été réalisée sur la qualité de l'eau dans une région, en prélevant des échantillons dans deux types de milieux : des rivières et des lacs. Chaque échantillon a ensuite été analysé et classé comme de qualité \og satisfaisante \fg{} ou \og insatisfaisante \fg{} selon les normes en vigueur.

Voici les données recueillies à partir d'un total de 500 échantillons analysés :
\begin{itemize}
	\item 300 échantillons provenaient de rivières ;
	\item parmi les échantillons prélevés dans les lacs, 80 étaient classés de qualité \og insatisfaisante \fg{} ;
	\item au total, 280 échantillons étaient classés de qualité \og satisfaisante \fg.
\end{itemize}

\begin{enumerate}
	\item Recopier et compléter le tableau suivant :

	\begin{center}
		\begin{tblr}{hlines, vlines, colspec={l*{3}{X[c,25mm]}}}
			& Rivière & Lac & Total \\
			Qualité satisfaisante & & & \\
			Qualité insatisfaisante & & & \\
			Total & & & 500 \\
		\end{tblr}
	\end{center}

	\item On choisit un échantillon d'eau au hasard parmi les 500 échantillons analysés.

	On note les évènements suivants :
	\begin{itemize}[left=5mm,label={}]
		\item $S$ : \og l'échantillon choisi est classé de qualité satisfaisante \fg{} ;
		\item $R$ : \og l'échantillon choisi provient d'une rivière \fg.
	\end{itemize}

	\emph{Pour un évènement quelconque $E$ : on note $\overline{E}$ l'évènement contraire de $E$ et $P(E)$ la probabilité de $E$.}

	On donnera les résultats sous forme de fraction.

	\begin{enumerate}
		\item Quelle est la probabilité que l'échantillon provienne d'une rivière ?
		\item Décrire l'évènement $\overline{S}$ par une phrase puis justifier que $P\left(\overline{S}\right) = \dfrac{11}{25}$.
		\item Calculer $P\left(S \cap \overline{R}\right)$, puis interpréter ce résultat dans le contexte de l'énoncé.
	\end{enumerate}

	\item La situation est représentée par l'arbre pondéré ci-dessous :

	\begin{center}
		\begin{tikzpicture}[x=3cm, y=1cm]
			
			\node (root) at (0,0) {};
			\node (A) at (1,1) {$R$};
			\node (Abar) at (1,-1) {$\overline{R}$};
			\node (B1) at (2,1.6) {$S$};
			\node (B1b) at (2,0.4) {$\overline{S}$};
			\node (B2) at (2,-0.4) {$S$};
			\node (B2b) at (2,-1.6) {$\overline{S}$};
			
			\draw (root.east) -- (A) node[pos=0.5,above] {\ldots};
			\draw (root.east) -- (Abar) node[pos=0.5,below] {\ldots};
			
			\draw (A.east) -- (B1) node[pos=0.5,above] {\ldots};
			\draw (A.east) -- (B1b) node[pos=0.5,below] {$x$};
			\draw (Abar.east) -- (B2) node[pos=0.5,above] {\ldots};
			\draw (Abar.east) -- (B2b) node[pos=0.5,below] {\ldots};
		\end{tikzpicture}
	\end{center}

	\begin{enumerate}
		\item Traduire la probabilité $x$ en utilisant les évènements $\overline{S}$ et $R$.
		\item Vérifier que $x = \dfrac{7}{15}$.
		\item Recopier et compléter entièrement l'arbre.
	\end{enumerate}

	\item Les évènements \og l'échantillon provient d'une rivière \fg{} et \og l'échantillon est de qualité insatisfaisante \fg{} sont-ils indépendants ?
\end{enumerate}

\end{document}
